% SEGMENT OLP-0670-S01
% Part: proof-theory
% Chapter: natural-deduction
% Section: translation-G2i

\documentclass[../../../include/open-logic-section]{subfiles}

\begin{document}

\olfileid{pt}{nat}{tgi}

\olsection{\Log{G2i} இலிருந்து \Log{N2i} க்கு மாற்றுதல்}

\Log{G2i} தொடரணிகளின் இடப்புறங்கள் !!{formula}s அடங்கிய
பன்மைக் கணங்கள்~$\Gamma$; ஆனால் \Log{N2i} தொடரணிகளின்
இடப்புறங்கள், ஒவ்வொரு அடையாளமும் ஒருமுறை மட்டுமே வரும்
அடையாளமிட்ட !!{formula}s இன் கணங்கள்~$\Gamma'$. ஒவ்வொரு
!!{formula}~$!A$ க்கும், $\Gamma'$ இல் $x : !A$ வடிவிலான
!!{element}s இன் எண்ணிக்கை, $\Gamma$ இல் $!A$ இன்
நிகழ்வெண்ணை, அதாவது $!A$ இன் பிரதிகளின் எண்ணிக்கையை,
விடக் கூடாமலிருந்தால், அடையாளமிட்ட வாய்ப்பாடுகளின்
கணம்~$\Gamma'$ பன்மைக் கணம்~$\Gamma$ க்கு
\emph{ஒத்திருக்கிறது} எனக் கூறுவோம்.

\begin{prop}\ollabel{prop:G2i-to-N2i}
$\Log{G2i} + \Cut \Proves \Gamma \Sequent \Delta$ எனில்,
$\Log{N2i} \Proves \Gamma' \Sequent !C$; இதில் $\Delta$
வெற்றாக இருந்தால் $!C \ident \lfalse$, இல்லையெனில்
$\Delta = \{!A\}$; அப்போது முடிவின் வலப்புறம் அந்த
ஒற்றை வாய்ப்பாடாகும். மேலும் $\Gamma'$ ஆனது~$\Gamma$ க்கு
ஒத்திருக்கிறது.
\end{prop}

\begin{proof}
  $\pi$ என்பது \Log{G2i} மற்றும் \Cut இல் $\Gamma \Sequent
  \Delta$ இன் !!a{proof} எனில், $\Gamma$ க்கு ஒத்த~$\Gamma'$
  ஒன்றும், \Log{N2i} இல் $\Gamma' \Sequent !C$ இன்
  !!a{proof}~$\delta$ ஒன்றும் உள்ளன என்று காட்டுவோம்.
  இதை $n = \pheight{\pi}$ மீது தொகுத்தறிந்து நிறுவுகிறோம்.

  $n = 0$ எனில், $\pi$ ஓர் அடிகோள். அது $\lfalse \Sequent
  \quad$ என்ற அடிகோளாக இருந்தால், $\delta$ என்பது
  $x:\lfalse \Sequent \lfalse$ என்ற அடிகோள்; அதாவது
  $\Gamma' = \{x:\lfalse\}$ மற்றும் $!C \ident \lfalse$.
  $\pi$ ஆனது $!A \Sequent !A$ வடிவிலான அடிகோள் எனில்,
  $\delta$ என்பது $x: !A \Sequent !A$.

% SEGMENT OLP-0670-S02
  $n>0$ எனில், $\pi$ ஒரு விதி~$R$ இல் முடியும். அந்த
  விதியின் முன்கோள்களுக்கு ஒத்த தொடரணிகளின் !!{proof}s
  \Log{N2i} இல் உள்ளன என்று தொகுத்தறிதல் கருதுகோள்
  உறுதி செய்கிறது. அவற்றில் விடுவிப்பு அடையாளங்கள் அனைத்தும்
  வெவ்வேறானவை என்றும், ஓர் அடையாளம்~$x$ விடுவிக்கப்படுமானால்
  அதை விடுவிக்கும் அனுமானத்திற்கு மேலே மட்டுமே அது வரும்
  என்றும், தேவைப்பட்டால் அடையாளங்களை மாற்றி அமைத்துக்
  கொள்ளலாம். இப்போது இந்த !!{proof}s ஐ இணைத்து ஏற்ற
  $\Gamma' \Sequent !C$ இன் நிறுவலை எவ்வாறு பெறுவது எனக்
  காண்போம். $R$ க்கு ஏற்ப நேர்வுகளைப் பிரிப்போம்.

  \begin{enumerate}
    \item $R$ என்பது \RightR{\Weakening} எனில்,
    $\pi$ பின்வருமாறு முடியும்:
    \[
    \AxiomC{}
    \RightLabel{$\pi_1$}
    \Deduce$\Gamma \fCenter$
    \RightLabel{\RightR{\Weakening}}
    \UnaryInf$\Gamma \fCenter !A$
    \DisplayProof
    \]
    $\pi_1$ க்குத் தொகுத்தறிதல் கருதுகோளைப் பயன்படுத்தினால்
    $\Gamma' \Sequent \lfalse$ இன் !!a{proof} கிடைக்கும்.
    அதற்கு $\FalseInt$ ஐப் பயன்படுத்தி $\Gamma' \Sequent !A$
    இன் !!a{proof} பெறுகிறோம்.

% SEGMENT OLP-0670-S03
    \item $R$ என்பது \LeftR{\Weakening} எனில்,
    $\pi$ பின்வருமாறு முடியும்:
    \[
    \AxiomC{}
    \RightLabel{$\pi_1$}
    \Deduce$\Gamma \fCenter \Delta$
    \RightLabel{\LeftR{\Weakening}}
    \UnaryInf$!A, \Gamma \fCenter \Delta$
    \DisplayProof
    \]
    $\pi_1$ க்குத் தொகுத்தறிதல் கருதுகோளைப் பயன்படுத்தினால்
    $\Gamma' \Sequent !C$ இன் !!a{proof} கிடைக்கும்.
    இதையே $\delta$ ஆக எடுத்துக்கொள்வோம். $\Gamma'$ ஆனது
    $\Gamma$ க்கு ஒத்திருந்தால், அது $!A, \Gamma$ க்கும்
    ஒத்திருக்கிறது: $\Gamma'$ இலுள்ள $x : !A$ களின்
    எண்ணிக்கை, $\Gamma$ இலுள்ள~$!A$ பிரதிகளின்
    எண்ணிக்கையை விட அதிகமல்ல; எனவே $\Gamma \cup \{!A\}$
    இலுள்ள $!A$ பிரதிகளின் எண்ணிக்கையையும் விட அதிகமல்ல.

% SEGMENT OLP-0670-S04
    \item $R$ என்பது \LeftR{\Contraction} எனில்,
    $\pi$ பின்வருமாறு முடியும்:
    \[
    \AxiomC{}
    \RightLabel{$\pi_1$}
    \Deduce$!A, !A, \Gamma \fCenter \Delta$
    \RightLabel{\LeftR{\Contraction}}
    \UnaryInf$!A, \Gamma \fCenter \Delta$
    \DisplayProof
    \]
    $\pi_1$ க்குத் தொகுத்தறிதல் கருதுகோளைப் பயன்படுத்தினால்,
    $\Pi \subseteq \{x : !A, y:!A\}$ ஆகும் வகையில்
    $\Pi, \Gamma' \Sequent !C$ இன் !!a{proof} கிடைக்கும்.
    $\Pi = \{x : !A, y:!A\}$ எனில், அப்படி கிடைத்த
    நிறுவலில் வரும் எல்லா அடையாளமிட்ட
    !!{formula}s~$y:!A$ ஐயும் $x:!A$ ஆக மாற்றுவோம்.
    $x$ அதன் இறுதித் தொடரணியின் சூழலில் வருவதால்,
    அந்த நிறுவலில் எந்த அனுமானமும் $x:!B$ வடிவிலான
    !!a{formula} ஐ விடுவிக்காது என்று எடுத்துக்கொள்ளலாம்;
    அதாவது அந்த நிறுவலின் தொடரணி இடப்புறத்தில் வரும்
    எந்த $x:!B$ உம் செயல்நிலையிலிருக்காது. ஆகவே மாற்றிய
    பிறகும் கிடைப்பது \Log{N2i} இன் !!a{proof} ஆகும்.

% SEGMENT OLP-0670-S05
    \item $R$ என்பது $\RightR{\lif}$ எனில்,
    $\pi$ பின்வருமாறு முடியும்:
    \[
    \AxiomC{}
    \RightLabel{$\pi_1$}
    \Deduce$!A, \Gamma \fCenter !B$
    \RightLabel{\RightR{\lif}}
    \UnaryInf$\Gamma \fCenter !A \lif !B$
    \DisplayProof
    \]
    தொகுத்தறிதல் கருதுகோளால், $x: !A, \Gamma'
    \Sequent !B$ அல்லது $\Gamma' \Sequent !B$ இன்
    !!a{proof}~$\delta_1$ உள்ளது; இங்கு $\Gamma'$ ஆனது
    $\Gamma$ க்கு ஒத்திருக்கிறது. $\delta_1$ இன்
    இறுதித் தொடரணிக்கு \Intro{\lif} விதியைப் பயன்படுத்தி
    $\delta$ ஐப் பெறலாம்.

% SEGMENT OLP-0670-S06
    \item $R$ என்பது \LeftR{\lif} எனில்,
    $\pi$ பின்வருமாறு முடியும்:
    \[
    \AxiomC{}
    \RightLabel{$\pi_1$}
    \Deduce$\Gamma_1 \fCenter !A$
    \AxiomC{}
    \RightLabel{$\pi_2$}
    \Deduce$!B, \Gamma_2 \fCenter \Delta$
    \RightLabel{\LeftR{\lif}}
    \BinaryInf$!A \lif !B, \Gamma_1, \Gamma_2 \fCenter \Delta$
    \DisplayProof
    \]
    தொகுத்தறிதல் கருதுகோளால் $\Gamma_1' \Sequent !A$
    இன் !!{proof}~$\delta_1$ உம், $x: !B, \Gamma_2'
    \Sequent !C$ அல்லது $\Gamma_2' \Sequent !C$ இன்
    நிறுவல்~$\delta_2$ உம் கிடைக்கின்றன. இரண்டாம் வடிவில்
    $\delta_2$ இருந்தால், அதையே~$\delta$ ஆகக் கொள்ளலாம்.
    இல்லையெனில், $\delta_1$ மற்றும்~$\delta_2$ இரண்டிலும்
    ஒரே அடையாளம் வராதவாறு $\delta_1$ இலுள்ள வாய்ப்பாட்டு
    அடையாளங்களையும் விடுவிப்பு அடையாளங்களையும் மாற்றலாம்.
    அப்போது $\Gamma_1' \cup \Gamma_2'$ ஆனது $\Gamma_1
    \cup \Gamma_2$ க்கு ஒத்திருக்கிறது. $x:!B$ ஆனது
    $\delta_2$ இன் இறுதித் தொடரணியில் வருவதால்,
    $\delta_2$ இல் விடுவிப்பு அடையாளம்~$x$ உடைய எந்த
    அனுமானத்திலும் $x:!B$ செயல்நிலையிலிருக்க முடியாது.
    $\delta_2$ இலுள்ள ஒவ்வொரு $x:!B \Sequent !B$
    அடிகோளுக்கும் பதிலாகப் பின்வரும் !!{proof}~$\delta_3$
    ஐப் பதிப்போம்:
    \[
    \Axiom$x: !A \lif !B \fCenter !A \lif !B$
    \AxiomC{}
    \RightLabel{$\delta_1$}
    \Deduce$\Gamma_1' \fCenter !A$
    \RightLabel{\Elim{\lif}}
    \BinaryInf$x: !A \lif !B, \Gamma_1' \fCenter !B$
    \DisplayProof
    \]
    மேலும் $\delta_2$ இலுள்ள ஒவ்வொரு $x:!B$ ஐயும்
    $x:!A \lif !B$ ஆக மாற்றுவோம். இவ்வாறு கிடைக்கும்
    $\Subst{\delta_2}{\delta_3}{x:!B}$ ஆனது
    $x:!A \lif !B, \Gamma_1', \Gamma_2' \Sequent !C$
    இன் !!a{proof} ஆகும்.

% SEGMENT OLP-0670-S07
    \item $R$ என்பது \RightR{\land} எனில்,
    $\pi$ பின்வருமாறு முடியும்:
    \[
    \AxiomC{}
    \RightLabel{$\pi_1$}
    \Deduce$\Gamma_1 \fCenter !A$
    \AxiomC{}
    \RightLabel{$\pi_2$}
    \Deduce$\Gamma_2 \fCenter !B$
    \RightLabel{\RightR{\land}}
    \BinaryInf$\Gamma_1, \Gamma_2 \fCenter !A \land !B$
    \DisplayProof
    \]
    தொகுத்தறிதல் கருதுகோளால் $\Gamma_1' \Sequent !A$
    இன் நிறுவல்~$\delta_1$ உம், $\Gamma_2' \Sequent !B$
    இன் நிறுவல்~$\delta_2$ உம் கிடைக்கும். $\delta_2$
    இலுள்ள எல்லா வாய்ப்பாட்டு அடையாளங்களையும் விடுவிப்பு
    அடையாளங்களையும் மாற்றுவதன் மூலம், $\delta_1$ மற்றும்
    $\delta_2$ இரண்டிலும் ஒரே அடையாளம் வராமல் செய்யலாம்.
    அப்போது $\Gamma_1' \cup \Gamma_2'$ ஆனது $\Gamma_1
    \cup \Gamma_2$ க்கு ஒத்திருக்கிறது. $\delta_1$ மற்றும்
    $\delta_2$ இன் இறுதித் தொடரணிகளுக்கு
    $\Intro{\land}$ ஐப் பயன்படுத்தி !!a{proof}~$\delta$
    பெறுகிறோம்.

% SEGMENT OLP-0670-S08
    \item $R$ என்பது $\LeftR{\land}$ எனில்,
    எடுத்துக்காட்டாக $\pi$ பின்வருமாறு முடியும்:
    \[
    \AxiomC{}
    \RightLabel{$\pi_1$}
    \Deduce$!A, \Gamma \fCenter \Delta$
    \RightLabel{\LeftR{\land}}
    \UnaryInf$!A \land !B, \Gamma \fCenter \Delta$
    \DisplayProof
    \]
    தொகுத்தறிதல் கருதுகோளால் $x: !A, \Gamma'
    \Sequent !C$ அல்லது $\Gamma' \Sequent !C$ இன்
    !!a{proof}~$\delta_1$ உள்ளது; இதில் $\Gamma'$ ஆனது
    $\Gamma$ க்கு ஒத்திருக்கிறது. இரண்டாம் நேர்வில்
    $\delta_1$ ஐயே~$\delta$ ஆகக் கொள்ளலாம். முதல்
    நேர்வில், $x:!A$ இறுதித் தொடரணியில் வருவதால்,
    $\delta_1$ இல் விடுவிப்பு அடையாளம்~$x$ உடைய எந்த
    அனுமானத்திலும் $x:!A$ செயல்நிலையிலிருக்காது. ஒவ்வொரு
    $x:!A \Sequent !A$ அடிகோளுக்கும் பதிலாகப் பின்வரும்
    !!{proof}~$\delta_2$ ஐப் பதிப்போம்:
    \[
    \Axiom$x: !A \land !B \fCenter !A \land !B$
    \RightLabel{\Elim{\land}}
    \UnaryInf$x: !A \land !B \fCenter !A$
    \DisplayProof
    \]
    மேலும் ஒவ்வொரு $x:!A$ ஐயும் $x:!A \land !B$ ஆக
    மாற்றுவோம்; அதாவது
    $\Subst{\delta_1}{\delta_2}{x:!A}$ ஐக் கொள்வோம்.
    இது $x:!A \land !B, \Gamma' \Sequent !C$ இன்
    !!a{proof} ஆகும்.

% SEGMENT OLP-0670-S09
    \item $R$ என்பது \Cut எனில், $\pi$ பின்வருமாறு
    முடியும்:
    \[
    \AxiomC{}
    \RightLabel{$\pi_1$}
    \Deduce$\Gamma_1 \fCenter !A$
    \AxiomC{}
    \RightLabel{$\pi_2$}
    \Deduce$!A, \Gamma_2 \fCenter \Delta$
    \RightLabel{\Cut}
    \BinaryInf$\Gamma_1, \Gamma_2 \fCenter \Delta$
    \DisplayProof
    \]
    தொகுத்தறிதல் கருதுகோளால் $\Gamma_1' \Sequent !A$
    இன் நிறுவல்~$\delta_1$ உம், $x:!A, \Gamma_2'
    \Sequent !C$ அல்லது $\Gamma_2' \Sequent !C$ இன்
    நிறுவல்~$\delta_2$ உம் கிடைக்கும். இரண்டாம் நேர்வில்
    $\delta$ ஐ~$\delta_2$ ஆகக் கொள்வோம். இல்லையெனில்,
    $\delta_2$ இலுள்ள ஒவ்வொரு $x:!A \Sequent !A$
    அடிகோளையும் $\delta_1$ ஆல் மாற்றி,
    $\Subst{\delta_2}{\delta_1}{x:!A}$ என்ற
    $\Gamma_1', \Gamma_2' \Sequent !C$ இன்
    !!a{proof}~$\delta$ ஐப் பெறுகிறோம்.
  \end{enumerate}
\end{proof}

பதிப்புக் குறிப்பு: தொகுத்தறிதலின் மற்ற விதி நேர்வுகள்
உறைந்த மூலத்தில் தரப்படவில்லை.

\end{document}
